\section{Introduction}

The databases of Cummins--Pauli \cite{CumminsPauli2003} and of the LMFDB
\cite{LMFDB,RSZB2022} list congruence subgroups of $\PSL{2}{\ZZ}$ and open
subgroups of $\GL{2}{\ProfiniteIntegers}$ respectively, together with the
modular curves they define. The aim of this document is to set up the theory
needed to extend these databases to quotients of modular curves, such as
$\ModCurve{0}(N)/\langle\ALmatrix{e}\rangle$ or $\ModCurve{\mathrm{ns}}^{+}(p)$
modulo extra automorphisms.

The guiding picture is Conway's ``big picture'' \cite{Conway1996}: the set
of homothety classes of lattices in $\QQ^{2}$, on which
$\RationalProjGroup$, the commensurator of $\PSL{2}{\ZZ}$, acts. Conway
used it to describe the groups $\GammaZero{N}$, their Atkin--Lehner
extensions and their normalisers, and to explain the role of the number
$24$. We work instead with lattices in $\FiniteAdeles{\QQ}^{2}$ and the group
\[
  \AdelicGroup = \GL{2}{\FiniteAdeles{\QQ}}/\PositiveRationals ,
\]
and deduce the classical statements by intersecting with
$\RationalLiftGroup$. The group $\AdelicGroup$ contains
$\GL{2}{\ProfiniteIntegers}$, so that the subgroups studied in the LMFDB
appear unchanged, and it contains elements of finite order that act as
Atkin--Lehner involutions, so that quotients by them appear as
(compact open) levels as well.

\subsection*{Overview}
\Cref{sec:conventions} fixes all conventions (column vectors, left actions,
level structures, the lattice attached to $\tau$) and compares them with
the literature. \Cref{sec:ambient} introduces the ambient groups and the
three versions $\adelic{\Gamma}$, $\lifted{\Gamma}$, $\Gamma$ of a group.
\Cref{sec:bigpicture} describes the adelic big picture and the standard
groups in it, and \cref{sec:trees} collects facts about fixed points in
trees. \Cref{sec:helling} proves an adelic version of Helling's theorem
\cite{Helling1970}, from which the classical one follows.
\Cref{sec:centring} defines centred groups and the Helling level,
\cref{sec:principal} the principal congruence subgroups
$\GammaPrincipal{N}{M}$ of the Atkin--Lehner groups, and
\cref{sec:quotients} describes the finite quotients they give rise to,
including a computational model. \Cref{sec:conjugacy} explains conjugacy
classes in fibre products, \cref{sec:automorphisms} modular automorphisms
and their Galois action, and \cref{sec:labels} summarises the state of the
labelling scheme.

\todo[inline]{Introduction to be expanded once the content has stabilised.}
